% =========================================================================
% ACESSO AO CÓDIGO-FONTE E REPOSITÓRIO GITHUB
% =========================================================================
\newpage
\apendice{Acesso ao Código-Fonte e Repositório GitHub}
\label{ap:github}

Para garantir a reprodutibilidade dos experimentos e disponibilizar o código-fonte desenvolvido, todo o material está organizado em repositório público no GitHub.

O repositório \textbf{Networks Flow} está estruturado nos seguintes diretórios principais:
\begin{itemize}[leftmargin=1.5em, itemsep=0.2em]
\raggedright
    \item \textbf{Fluxo Máximo} (\path{Implementações/FlowNetwork/}): Módulos base e algoritmos (\lstinline{FordFulkerson}, \lstinline{EdmondsKarp}, \lstinline{Dinic}, \lstinline{PushRelabel} e \lstinline{PushRelabelImproved});
    \item \textbf{Fluxo de Custo Mínimo} (\path{Implementações/CostNetwork/}): Módulos base e algoritmos (\lstinline{CycleCanceling}, \lstinline{SuccessiveShortestPath} e \lstinline{NetworkSimplex});
    \item \textbf{Problemas} (\path{Implementações/Problemas/}): Resoluções para problemas combinatórios clássicos (como \textit{Minimum Vertex Cover}, emparelhamentos e caminhos disjuntos);
    \item \textbf{Aplicações Práticas} (\path{Aplicações/Segmentação de Imagens/}): Segmentação de imagens via corte mínimo em grafos (\textit{Graph Cuts});
    \item \textbf{Experimentos} (\path{Experimentos/}): Scripts de automação de testes, leitura de instâncias DIMACS e avaliação experimental.
\end{itemize}

O código-fonte completo pode ser consultado, clonado e executado através do endereço:

\vspace{0.4cm}
\begin{center}
    \fbox{\url{https://github.com/GabrielFrigo4/networks-flow}}
\end{center}
\vspace{0.4cm}

O repositório contém o histórico de desenvolvimento e as instruções para compilação e reprodução dos experimentos.
